\section{Conclusion, Roadmap and Availability}
\label{sec:conclusion}

Language-model agents can now hold money and hand it to the agents they hire; until they are measured against the
best split that money could buy, how well they spend it is a matter of belief. \sys{} v2 makes agents allocate, with
reserved pools in dollars and memory, refunds on firing, and make or buy over open protocols, and records every
allocation. \bench{} makes the optimum known and holds itself to measuring allocation before any allocator runs. On
it, agents that hold the money first spent it worse than every rule we compared them with, and the traces showed
why; fixing the money mechanics closed half of the gap, and the half that remains is the agents' own judgment.

\textbf{Roadmap.} The remaining gap points at the agents rather than the mechanics. The CEO never used \texttt{attach},
so handing work over whole should not depend on the agent remembering to: a hirer could pass a job by reference by
default. The CEO spends over half its budget rereading the request on every turn; keeping long inputs out of its
context, or caching them, would leave that money for the work. Declared difficulty barely tracks real difficulty, so
the router could learn from outcomes instead of trusting the declaration. The next measurements are the ones these runs
left out: a recalibration with harder top levels and hidden bins that change success rates (checks V1, V4 and V5),
more projects and seeds, local models (and with them memory as a second currency), a small cloud tier, outside agents
whose reliability is unknown, and the cost in time of the protocol boundaries.

\textbf{Availability.} \sys{} v2 and \bench{} are source-available under the PolyForm Noncommercial License 1.0.0 at
\url{https://github.com/HASHIRU-AI/HASHIRU}, with the outcome table, every run's results and traces, and the
environment of each run. \sys{} v1's code is preserved at tag \texttt{v1.0-paper}.
